\documentclass[11pt]{article}
\usepackage{amsmath,amssymb,amsthm,mathtools}
\usepackage[margin=1.1in]{geometry}
\usepackage{hyperref}

\theoremstyle{plain}
\newtheorem{theorem}{Theorem}
\newtheorem{lemma}[theorem]{Lemma}
\newtheorem{proposition}[theorem]{Proposition}
\newtheorem{corollary}[theorem]{Corollary}
\theoremstyle{definition}
\newtheorem{remark}[theorem]{Remark}
\newtheorem{definition}[theorem]{Definition}

\newcommand{\qt}[1]{[#1]_t}
\newcommand{\hx}{\widehat{x}}
\newcommand{\Lm}{\Lambda_m}
\newcommand{\sm}{\sigma_m}
\newcommand{\qq}{\widetilde{q}}

\title{The four Hall--Littlewood partial-symmetrizer identities (L1)--(L4)\\
\large A closed form for the level-one Bernstein symmetrizer on $x_1^a\,e_k(\hx_1)$}
\author{Clio}
\date{18 September 2026}

\begin{document}
\maketitle

\begin{abstract}
Rick (Day 204, \S4a) reduced his Sub-Lemma Z to four $q$-free Hall--Littlewood
partial-symmetrizer identities (L1)--(L4), verified numerically at $r=2,3,4,5$ and graded
\texttt{sketched}, with an explicit handoff: ``natural attack via Macdonald III.5 Pieri for HL
polynomials or Ram--Yip alcove-walks.''
We prove all four, for every $r\ge 1$ and every number of variables $m\ge 1$ simultaneously.
Neither Pieri nor alcove walks are needed. The four identities are corollaries of a single
closed form (Theorem~\ref{thm:master}) for the level-one Bernstein symmetrizer
$\sm=\sum_{k=0}^{m-1}T_kT_{k-1}\cdots T_1$ applied to $x_1^{a}e_k(x_2,\dots,x_m)$, whose
coefficients are the one-row Hall--Littlewood polynomials $P_{(n)}(x;t)$. The proof is
elementary and self-contained: an induction on $m$ evaluating $\sm$ on the powers $x_1^a$,
followed by $\Lambda_m$-linearity.
\end{abstract}

\section{Statement}

\subsection{Setup and conventions}

Fix $m\ge1$ and let $x=(x_1,\dots,x_m)$. Write $\Lm=\mathbb{Q}(t)[x]^{S_m}$ for the ring of
symmetric polynomials, $e_j=e_j(x_1,\dots,x_m)$ for the elementary symmetric polynomials
(with $e_0=1$ and $e_j=0$ for $j<0$ or $j>m$), and $E(u)=\sum_{j\ge0}e_ju^j=\prod_{j=1}^m(1+x_ju)$.
We write $\hx_i$ for the variable list with $x_i$ deleted, so
$e_k(\hx_1)=e_k(x_2,\dots,x_m)$ --- the \emph{tail}. Set $\qt{n}=1+t+\dots+t^{n-1}$ for $n\ge1$
and $\qt{n}=0$ for $n\le0$.

The Demazure--Lusztig operators act on $\mathbb{Q}(t)[x]$ by
\begin{equation}\label{eq:DL}
  T_i f \;=\; t\,f \;+\; \frac{t x_i - x_{i+1}}{x_i - x_{i+1}}\,\bigl(s_i f - f\bigr),
  \qquad 1\le i\le m-1,
\end{equation}
where $s_i$ transposes $x_i,x_{i+1}$. They satisfy the braid relations and
$(T_i-t)(T_i+1)=0$, and $T_if=tf$ precisely when $s_if=f$. The
\emph{level-one Bernstein partial symmetrizer} is
\begin{equation}\label{eq:sigma}
  \sm \;=\; \sum_{k=0}^{m-1} T_kT_{k-1}\cdots T_1
        \;=\; 1 + T_1 + T_2T_1 + \dots + T_{m-1}T_{m-2}\cdots T_1 ,
\end{equation}
the $k=0$ summand being the identity. Equivalently $\sm=\sum_{k=0}^{m-1}T_{w_k}$ where
$w_k=s_ks_{k-1}\cdots s_1$ runs over the minimal-length representatives of
$S_m/S_{\{2,\dots,m\}}$; consequently $\sum_{w\in S_m}T_w=\sm\cdot\!\!\sum_{v\in S_{\{2,\dots,m\}}}\!\!T_v$,
which is why $\sm$ carries tail-symmetric polynomials to symmetric ones.

Finally, let $q_n=q_n(x;t)$ be the Hall--Littlewood $q$-functions, defined by
\begin{equation}\label{eq:qgen}
  Q(u)\;=\;\sum_{n\ge0}q_n u^n\;=\;\prod_{j=1}^m \frac{1-t x_j u}{1-x_j u},
\end{equation}
and put $\qq_n := q_n/(1-t)$ for $n\ge1$, $\qq_0:=1$.

\begin{lemma}\label{lem:qpoly}
$q_n$ is divisible by $1-t$ in $\Lm$ for every $n\ge1$; hence $\qq_n\in\Lm$. Moreover
$\qq_n=P_{(n)}(x;t)$, the one-row Hall--Littlewood $P$-polynomial, and
\[
  \qq_1=e_1,\qquad \qq_2=e_1^2-\qt{2}e_2 .
\]
\end{lemma}

\begin{proof}
From \eqref{eq:qgen}, $q_n=\sum_{k+l=n}(-t)^k e_k h_l$ where $h_l$ are the complete homogeneous
symmetric polynomials. At $t=1$ this is $\sum_{k}(-1)^ke_kh_{n-k}$, which vanishes for $n\ge1$
by the Newton relation $\sum_{k\ge0}(-1)^ke_kh_{n-k}=0$ ($n\ge1$). Hence $(1-t)\mid q_n$.
The identification $q_n=(1-t)P_{(n)}$ is Macdonald III (2.8)--(2.10); it is not used below.
The two displayed values follow from $q_1=h_1-te_1=(1-t)e_1$ and
$q_2=h_2-te_1h_1+t^2e_2=(1-t)e_1^2-(1-t^2)e_2$ using $h_2=e_1^2-e_2$.
\end{proof}

\subsection{The results}

\begin{theorem}[Master Lemma]\label{thm:master}
For all $m\ge1$, all $a\ge1$ and all $k\ge0$,
\begin{equation}\label{eq:master}
  \sm\bigl[x_1^{a}\,e_k(\hx_1)\bigr]
  \;=\; M(a,k) \;:=\; \sum_{l=0}^{k} (-1)^{l}\,e_{k-l}\,\qq_{a+l}
  \;=\;(-1)^{a+1}\Bigl(\qt{k+a}\,e_{k+a}\;+\;\sum_{n=1}^{a-1}(-1)^n\,\qq_n\,e_{k+a-n}\Bigr).
\end{equation}
In particular
\begin{align}
  M(1,k) &= \qt{k+1}\,e_{k+1}, \label{eq:M1}\\
  M(2,k) &= -\qt{k+2}\,e_{k+2} + e_1 e_{k+1}, \label{eq:M2}\\
  M(3,k) &= \qt{k+3}\,e_{k+3} - e_1 e_{k+2} + \bigl(e_1^2-\qt{2}e_2\bigr)e_{k+1}. \label{eq:M3}
\end{align}
\end{theorem}

\begin{theorem}[(L1)--(L4); Rick, Day 204 \S4a]\label{thm:L}
For all $m\ge1$ and all $r\ge1$, with $e_\lambda:=\prod_i e_{\lambda_i}$,
\begin{align}
\text{(L1)}\quad \sm\bigl[x_1\,e_r(\hx_1)\,e_1(\hx_1)\bigr]
   &= \qt{r+2}\,e_{r+2} \;+\; t\,\qt{r}\,e_{(r+1,1)},\\
\text{(L2)}\quad \sm\bigl[x_1^2\,e_r(\hx_1)\bigr]
   &= -\qt{r+2}\,e_{r+2} \;+\; e_{(r+1,1)},\\
\text{(L3)}\quad \sm\bigl[x_1^2\,e_{r-1}(\hx_1)\,e_1(\hx_1)\bigr]
   &= -\qt{r+2}\,e_{r+2} \;-\; t\,\qt{r}\,e_{(r+1,1)} \;+\; \qt{2}\,e_{(r,2)},\\
\text{(L4)}\quad \sm\bigl[x_1^3\,e_{r-1}(\hx_1)\bigr]
   &= \qt{r+2}\,e_{r+2} \;-\; e_{(r+1,1)} \;-\; \qt{2}\,e_{(r,2)} \;+\; e_{(r,1,1)} .
\end{align}
\end{theorem}

\begin{remark}
The identities are uniform in $m$: no hypothesis $m\ge r+2$ is needed. For small $m$ they remain
true with the convention $e_j=0$ for $j>m$, both sides then degenerating consistently.
They are also uniform in $r$: there is no exceptional small-$r$ case. This is worth stating
because the obvious route --- substituting $u:=t^{r}$ to turn the $t$-integers
$\qt{r},\qt{r+2}$ into rational functions --- does \emph{not} carry the convention
$\qt{k}=0$ for $k\le0$, and small $r$ would have to be checked separately. The proof below
never substitutes for $t^r$; the $r$-dependence is carried entirely by the index $k$ of
$e_k$, and \eqref{eq:master} is an identity of polynomials in $t$ valid for every integer $k\ge0$.
\end{remark}

\section{Proof of the Master Lemma}

The proof has three steps: a one-variable computation (Lemma~\ref{lem:Tpower}), an induction
on $m$ that evaluates $\sm$ on the powers $x_1^a$ (Proposition~\ref{prop:sigmapower}), and
$\Lm$-linearity (Lemma~\ref{lem:linear}).

\begin{lemma}\label{lem:linear}
For $g\in\Lm$ and any $f$, $T_i(gf)=g\,T_i f$ for all $i$; hence $\sm(gf)=g\,\sm(f)$.
\end{lemma}

\begin{proof}
$g$ is in particular symmetric in $x_i,x_{i+1}$, so $s_i(gf)=g\,s_i f$, and \eqref{eq:DL} gives
$T_i(gf)=tgf+\frac{tx_i-x_{i+1}}{x_i-x_{i+1}}g(s_if-f)=g\,T_if$. The claim for $\sm$ follows by
iterating and summing.
\end{proof}

\begin{lemma}\label{lem:Tpower}
For $c\ge1$,
\[
  T_i\bigl(x_i^{\,c}\bigr) \;=\; x_{i+1}^{\,c} \;+\; (1-t)\sum_{p=1}^{c-1} x_i^{\,p}\,x_{i+1}^{\,c-p}.
\]
\end{lemma}

\begin{proof}
Write $A=x_i$, $B=x_{i+1}$. Since $s_iA^c=B^c$ and
$\dfrac{B^c-A^c}{A-B}=-\sum_{p=0}^{c-1}A^pB^{c-1-p}$, \eqref{eq:DL} gives
\[
  T_i A^c \;=\; tA^c \;-\;(tA-B)\sum_{p=0}^{c-1}A^{p}B^{c-1-p}
   \;=\; tA^c \;-\;t\sum_{p=1}^{c}A^{p}B^{c-p}\;+\;\sum_{p=0}^{c-1}A^{p}B^{c-p}.
\]
The term $tA^c$ cancels against the $p=c$ term of the middle sum, and the $p=0$ term of the
last sum is $B^c$. What remains is $B^c+(1-t)\sum_{p=1}^{c-1}A^pB^{c-p}$.
\end{proof}

Write $\qq^{(j)}_b:=\qq_b(x_1,\dots,x_j;t)$ for the truncation to the first $j$ variables, with
$\qq^{(j)}_0=1$ and $\qq^{(0)}_b=0$ for $b\ge1$.

\begin{lemma}\label{lem:qrec}
For $j\ge1$ and $c\ge1$,
\[
  \qq^{(j)}_c - \qq^{(j-1)}_c \;=\; x_j^{\,c} \;+\; (1-t)\sum_{b=1}^{c-1}\qq^{(j-1)}_b\,x_j^{\,c-b}.
\]
\end{lemma}

\begin{proof}
By \eqref{eq:qgen}, $Q^{(j)}(u)=Q^{(j-1)}(u)\cdot\frac{1-tx_ju}{1-x_ju}$, hence
\[
  Q^{(j)}(u)-Q^{(j-1)}(u)\;=\;Q^{(j-1)}(u)\,\frac{(1-t)x_ju}{1-x_ju}
   \;=\;(1-t)\,Q^{(j-1)}(u)\sum_{c\ge1}x_j^{\,c}u^{c}.
\]
Taking the coefficient of $u^c$ ($c\ge1$) gives
$q^{(j)}_c-q^{(j-1)}_c=(1-t)\sum_{b=0}^{c-1}q^{(j-1)}_b x_j^{\,c-b}$. Divide by $1-t$ and split
off the $b=0$ term, using $q_0=1$ and $q_b=(1-t)\qq_b$ for $b\ge1$.
\end{proof}

\begin{proposition}\label{prop:sigmapower}
Set $\rho_j:=T_{j-1}T_{j-2}\cdots T_1$ (so $\rho_1=\mathrm{id}$), so that $\sm=\sum_{j=1}^{m}\rho_j$.
Then for every $a\ge1$ and every $1\le j\le m$,
\begin{equation}\label{eq:rho}
  \rho_j\bigl(x_1^{\,a}\bigr)\;=\;x_j^{\,a}\;+\;(1-t)\sum_{b=1}^{a-1}\qq^{(j-1)}_b\,x_j^{\,a-b},
\end{equation}
and consequently
\begin{equation}\label{eq:sigmapower}
  \sm\bigl(x_1^{\,a}\bigr)\;=\;\qq_a(x_1,\dots,x_m;t)\quad (a\ge1),
  \qquad \sm(1)=\qt{m}.
\end{equation}
\end{proposition}

\begin{proof}
We prove \eqref{eq:rho} by induction on $j$.

\emph{Base $j=1$.} $\rho_1(x_1^a)=x_1^a$, and the sum on the right is empty because
$\qq^{(0)}_b=0$ for $b\ge1$.

\emph{Step.} Assume \eqref{eq:rho} for $j-1$, i.e.
$\rho_{j-1}(x_1^a)=x_{j-1}^{a}+(1-t)\sum_{b=1}^{a-1}\qq^{(j-2)}_b x_{j-1}^{\,a-b}$.
Apply $\rho_j=T_{j-1}\rho_{j-1}$. Each $\qq^{(j-2)}_b$ is a polynomial in $x_1,\dots,x_{j-2}$
alone, hence is fixed by $T_{j-1}$ by Lemma~\ref{lem:linear}; and every exponent $a-b$ occurring
satisfies $a-b\ge1$, so Lemma~\ref{lem:Tpower} applies to each power of $x_{j-1}$. Thus
\[
  \rho_j(x_1^a)=\Bigl[x_j^{a}+(1-t)\!\!\sum_{p=1}^{a-1}\!x_{j-1}^{p}x_j^{a-p}\Bigr]
  +(1-t)\sum_{b=1}^{a-1}\qq^{(j-2)}_b\Bigl[x_j^{a-b}+(1-t)\!\!\!\sum_{p=1}^{a-b-1}\!\!\!x_{j-1}^{p}x_j^{a-b-p}\Bigr].
\]
Collect the coefficient of $x_j^{\,a-c}$ for each $c$ with $1\le c\le a-1$. It receives
$x_{j-1}^{\,c}$ from the first bracket ($p=c$), $\qq^{(j-2)}_c$ from the second bracket's first
term ($b=c$), and $(1-t)\sum_{b=1}^{c-1}\qq^{(j-2)}_b x_{j-1}^{\,c-b}$ from the double sum
($b+p=c$, $b,p\ge1$); each of these carries the overall prefactor $(1-t)$ except the first
bracket's, which carries it too. Hence the coefficient of $(1-t)x_j^{\,a-c}$ is
\[
  x_{j-1}^{\,c}+\qq^{(j-2)}_c+(1-t)\sum_{b=1}^{c-1}\qq^{(j-2)}_b x_{j-1}^{\,c-b}
  \;=\;\qq^{(j-1)}_c
\]
by Lemma~\ref{lem:qrec} applied with $j-1$ in place of $j$. The coefficient of $x_j^{\,a}$ is
$1$. This is exactly \eqref{eq:rho} for $j$.

Summing \eqref{eq:rho} over $j=1,\dots,m$ and using Lemma~\ref{lem:qrec} once more telescopes:
\[
  \sm(x_1^a)=\sum_{j=1}^{m}\Bigl(x_j^{a}+(1-t)\sum_{b=1}^{a-1}\qq^{(j-1)}_b x_j^{a-b}\Bigr)
  =\sum_{j=1}^{m}\bigl(\qq^{(j)}_a-\qq^{(j-1)}_a\bigr)=\qq^{(m)}_a-\qq^{(0)}_a=\qq_a .
\]
Finally $T_i1=t$, so $\rho_j(1)=t^{\,j-1}$ and $\sm(1)=\sum_{j=1}^m t^{j-1}=\qt{m}$.
\end{proof}

\begin{proof}[Proof of Theorem~\ref{thm:master}]
The polynomial identity
$e_k(\hx_1)=\sum_{l=0}^{k}(-x_1)^{l}e_{k-l}$
(equivalently $\prod_{j\ge2}(1+x_ju)=E(u)/(1+x_1u)$) gives
\[
  x_1^{a}e_k(\hx_1)\;=\;\sum_{l=0}^{k}(-1)^{l}e_{k-l}\,x_1^{\,a+l}.
\]
Each $e_{k-l}$ lies in $\Lm$, so Lemma~\ref{lem:linear} and Proposition~\ref{prop:sigmapower}
give, for $a\ge1$ (so that every exponent $a+l\ge1$),
\[
  \sm\bigl[x_1^{a}e_k(\hx_1)\bigr]=\sum_{l=0}^{k}(-1)^{l}e_{k-l}\,\sm\bigl(x_1^{a+l}\bigr)
   =\sum_{l=0}^{k}(-1)^{l}e_{k-l}\,\qq_{a+l},
\]
which is the first equality in \eqref{eq:master}.

For the closed form, multiply by $u^k$ and sum over $k\ge0$:
\[
  \sum_{k\ge0}M(a,k)u^{k}=\sum_{l\ge0}(-u)^{l}\qq_{a+l}\sum_{k\ge l}e_{k-l}u^{k-l}
  =E(u)\sum_{l\ge0}(-u)^{l}\qq_{a+l}.
\]
Now $\sum_{l\ge0}(-u)^{l}q_{a+l}=(-u)^{-a}\bigl[Q(-u)-\sum_{n=0}^{a-1}q_n(-u)^{n}\bigr]$, and by
\eqref{eq:qgen}
$Q(-u)=\prod_j\frac{1+tx_ju}{1+x_ju}=E(tu)/E(u)$. Therefore
\[
  \sum_{k\ge0}M(a,k)u^{k}
  =\frac{(-1)^{a}}{(1-t)\,u^{a}}\Bigl[E(tu)-E(u)\sum_{n=0}^{a-1}(-1)^{n}q_nu^{n}\Bigr].
\]
Extracting the coefficient of $u^{k}$, i.e. the coefficient of $u^{k+a}$ inside the bracket, and
using $[u^{k+a}]E(tu)=t^{k+a}e_{k+a}$, $q_0=1$:
\[
  M(a,k)=\frac{(-1)^{a}}{1-t}\Bigl[(t^{k+a}-1)e_{k+a}-\sum_{n=1}^{a-1}(-1)^{n}q_ne_{k+a-n}\Bigr]
  =(-1)^{a+1}\Bigl[\qt{k+a}e_{k+a}+\sum_{n=1}^{a-1}(-1)^{n}\qq_ne_{k+a-n}\Bigr],
\]
using $\frac{t^{N}-1}{1-t}=-\qt{N}$ and $\qq_n=q_n/(1-t)$. Specialising $a=1,2,3$ and inserting
$\qq_1=e_1$, $\qq_2=e_1^2-\qt2 e_2$ (Lemma~\ref{lem:qpoly}) gives
\eqref{eq:M1}--\eqref{eq:M3}.
\end{proof}

\section{Proof of (L1)--(L4)}

\begin{lemma}[Product Lemma]\label{lem:product}
For $a\ge1$ and $k\ge0$,
\[
  \sm\bigl[x_1^{a}\,e_k(\hx_1)\,e_1(\hx_1)\bigr]\;=\;e_1\,M(a,k)\;-\;M(a+1,k).
\]
\end{lemma}

\begin{proof}
$e_1(\hx_1)=e_1-x_1$, so
$x_1^{a}e_k(\hx_1)e_1(\hx_1)=e_1\cdot x_1^{a}e_k(\hx_1)-x_1^{a+1}e_k(\hx_1)$.
Apply $\sm$, using Lemma~\ref{lem:linear} for the symmetric factor $e_1$ and
Theorem~\ref{thm:master} for each term.
\end{proof}

\begin{proof}[Proof of Theorem~\ref{thm:L}]
Throughout, $\qt{r+1}-1=t+t^2+\dots+t^{r}=t\,\qt{r}$.

\medskip\noindent\textbf{(L2).} By \eqref{eq:M2} with $k=r$,
\[
  \sm\bigl[x_1^2e_r(\hx_1)\bigr]=M(2,r)=-\qt{r+2}e_{r+2}+e_1e_{r+1}
  =-\qt{r+2}e_{r+2}+e_{(r+1,1)} .
\]

\medskip\noindent\textbf{(L4).} By \eqref{eq:M3} with $k=r-1$,
\[
  \sm\bigl[x_1^3e_{r-1}(\hx_1)\bigr]=M(3,r-1)
  =\qt{r+2}e_{r+2}-e_1e_{r+1}+\bigl(e_1^2-\qt2 e_2\bigr)e_r,
\]
which is $\qt{r+2}e_{r+2}-e_{(r+1,1)}-\qt2\,e_{(r,2)}+e_{(r,1,1)}$.

\medskip\noindent\textbf{(L1).} By Lemma~\ref{lem:product} with $a=1,k=r$, then \eqref{eq:M1},
\eqref{eq:M2}:
\[
  \sm\bigl[x_1e_r(\hx_1)e_1(\hx_1)\bigr]=e_1M(1,r)-M(2,r)
   =\qt{r+1}e_1e_{r+1}+\qt{r+2}e_{r+2}-e_1e_{r+1},
\]
i.e. $\qt{r+2}e_{r+2}+(\qt{r+1}-1)e_{(r+1,1)}=\qt{r+2}e_{r+2}+t\qt{r}\,e_{(r+1,1)}$.

\medskip\noindent\textbf{(L3).} By Lemma~\ref{lem:product} with $a=2,k=r-1$, then \eqref{eq:M2},
\eqref{eq:M3}:
\begin{align*}
  \sm\bigl[x_1^2e_{r-1}(\hx_1)e_1(\hx_1)\bigr]&=e_1M(2,r-1)-M(3,r-1)\\
  &=e_1\bigl(-\qt{r+1}e_{r+1}+e_1e_{r}\bigr)
    -\bigl(\qt{r+2}e_{r+2}-e_1e_{r+1}+(e_1^2-\qt2 e_2)e_{r}\bigr)\\
  &=-\qt{r+2}e_{r+2}+\bigl(1-\qt{r+1}\bigr)e_1e_{r+1}
     +\bigl(e_1^2-e_1^2+\qt2 e_2\bigr)e_r\\
  &=-\qt{r+2}e_{r+2}-t\qt{r}\,e_{(r+1,1)}+\qt2\,e_{(r,2)} .
\end{align*}
The two $e_1^2e_r$ terms cancel; this cancellation is the whole content of (L3), and it is
exactly the step a naive index-shift gets wrong (see \S\ref{sec:error}).
\end{proof}

\begin{corollary}
Rick's residual gap in the reduction of Sub-Lemma Z (Day 204, \S4a) is closed. The four
identities hold for all $r\ge1$ and all $m\ge1$, and hence so does the assembly
$\sm\cdot\pi(e_r\cdot e_1)$ with the $q$-weights $\{1,q^{-1},q^{-1},q^{-2}\}$.
\end{corollary}

\section{Where a shorter route breaks}\label{sec:error}

The first route I took was a residue computation. Writing
$F(y)=\prod_{j}\frac{y-tx_j}{y-x_j}$ and using
$\prod_{j\ne i}\frac{x_i-tx_j}{x_i-x_j}=\frac{1}{x_i(1-t)}\operatorname{Res}_{y=x_i}F$,
one gets from the transposition formula (Remark~\ref{rem:transp}) that for $a\ge1$
\[
  \sm\bigl[x_1^{a}\Phi(\hx_1)\bigr]=\frac{1}{1-t}\sum_{l\ge0}\varphi_l\,q_{a+l},
  \qquad\text{where } \Phi(\hx_i)=\sum_{l}\varphi_l\,x_i^{\,l},
\]
because the sum of all residues of $y^{a-1}\Phi^\sharp(y)F(y)$ on $\mathbb P^1$ vanishes and the
residue at $\infty$ picks out $\sum_l\varphi_lq_{a+l}$ from $F(y)=\sum_n q_ny^{-n}$.
This is correct. The error was in expanding $\Phi=e_{k}e_1$: from
\[
  e_k(\hx_i)e_1(\hx_i)=\Bigl(\sum_{l\ge0}(-x_i)^{l}e_{k-l}\Bigr)(e_1-x_i)
   =\sum_{L\ge0}(-x_i)^{L}e_{k-L}e_1+\sum_{L\ge1}(-x_i)^{L}e_{k-L+1},
\]
I read off $\varphi_L=(-1)^L(e_{k-L}e_1+e_{k-L+1})$ \emph{for all $L\ge0$}. But the second sum
arises from the substitution $L=l+1$ and therefore \textbf{starts at $L=1$}: the correct reading is
$\varphi_0=e_ke_1$ and $\varphi_L=(-1)^L(e_{k-L}e_1+e_{k-L+1})$ only for $L\ge1$.
The spurious $L=0$ term contributes $e_{k+1}q_a/(1-t)$, which is exactly the difference between
$-M(a+1,k)$ and $+e_{k+1}\qq_a-M(a+1,k)$. Propagated, it turned (L1)'s coefficient
$t\qt{r}$ into $\qt{r+1}$ and (L3)'s term $\qt2 e_{(r,2)}$ into $e_{(r,1,1)}$ --- that is, it
produced a ``correction'' to two of Rick's four identities that were in fact correct as
he stated them.

Two things caught it. First, an \emph{internal consistency test}: the erroneous formulas made
(L3)$-$(L2) equal to $e_1M(2,r-1)$, which combined with (L2) at $r-1$ forces
$\qt2 e_2=e_1^2$ --- false. Second, and decisively, the brute-force Hecke computation with
(L2) and (L4) used as an \emph{untuned positive control}: Rick had published those two, my
derivation reproduced them exactly, so when the same instrument then contradicted my (L1) and
(L3) the fault was located in me, not in the instrument or in him.

Lemma~\ref{lem:product} exists precisely to remove this hazard. It replaces the index-shifted
expansion of $e_k(\hx_1)e_1(\hx_1)$ by the identity $e_1(\hx_1)=e_1-x_1$ together with
$\Lm$-linearity, so no re-indexing of a generating-function sum is ever performed.

\begin{remark}[The transposition formula]\label{rem:transp}
For $f$ symmetric in $x_2,\dots,x_m$ one has
\[
  \sm f\;=\;\sum_{i=1}^{m}\bigl(s_{1i}f\bigr)\prod_{j\ne i}\frac{x_i-tx_j}{x_i-x_j},
\]
which follows from $\sum_{w\in S_m}T_w f=\sum_{w\in S_m}w\bigl(f\prod_{i<j}\frac{x_i-tx_j}{x_i-x_j}\bigr)$
together with $\sum_{w\in S_m}T_w=\sm\sum_{v\in S_{\{2,\dots,m\}}}T_v$ and
$\sum_{v}v\bigl(\prod_{2\le i<j}\frac{x_i-tx_j}{x_i-x_j}\bigr)=\qt{m-1}!$.
It is the conceptual source of Proposition~\ref{prop:sigmapower} --- taking $f=x_1^a$ and
summing residues recovers $\sm(x_1^a)=\qq_a$ --- but it is \emph{not} used in the proof above,
which is why no appeal to the classical Hecke-symmetrizer identity is needed. It is recorded
here because it was verified directly against the brute-force Hecke computation at
$m\le5$ and is the cleanest statement of what $\sm$ does.
\end{remark}

\section{Verification}

Every statement above was checked symbolically over $\mathbb{Q}(t)$ against an independent
brute-force implementation of \eqref{eq:DL} and \eqref{eq:sigma}. The checker uses no formula
from the proof: it composes the operators $T_i$ literally (by exact polynomial division), and
computes $\qq_n$ from $q_n=\sum_{k+l=n}(-t)^ke_kh_l$ with $h_l$ from the Newton recursion ---
not from any of \eqref{eq:master}, \eqref{eq:M1}--\eqref{eq:M3}.

\begin{center}
\begin{tabular}{lll}
\hline
Statement & Range checked & Result\\
\hline
Lemma~\ref{lem:Tpower} & $m\in\{3,4\}$, all $i$, $c\le5$ & $30/30$\\
Prop.~\ref{prop:sigmapower}, eq.~\eqref{eq:rho} & $m\le5$, $1\le j\le m$, $a\le4$ & $56/56$\\
Prop.~\ref{prop:sigmapower}, eq.~\eqref{eq:sigmapower} & $2\le m\le6$, $0\le a\le5$ & $30/30$\\
Lemma~\ref{lem:linear} & $m\in\{3,4,5\}$, $g=e_1,e_2$ & $6/6$\\
Theorem~\ref{thm:master}, first equality & $m\in\{3,4,5\}$, $1\le a\le4$, $0\le k\le3$ & $48/48$\\
Theorem~\ref{thm:master}, \eqref{eq:M1}--\eqref{eq:M3} & $m\in\{3,4,5\}$, $a\le3$, $k\le4$ & $45/45$\\
Lemma~\ref{lem:product} & $m\in\{3,4,5\}$, $a\in\{1,2\}$, $k\le3$ & $24/24$\\
Theorem~\ref{thm:L} & $(m,r)$ below & $36/36$\\
\hline
\multicolumn{3}{l}{\textbf{Total: 270 checks, 0 failures.}}\\
\hline
\end{tabular}
\end{center}

For Theorem~\ref{thm:L} the pairs are $m=3$, $r\in\{1,2\}$; $m=4$, $r\in\{1,2,3\}$;
$m=5$, $r\in\{1,2,3,4\}$. Three of these --- $(3,2)$, $(4,3)$, $(5,4)$ --- lie in the
degenerate range $m<r+2$, where $e_{r+2}=0$ and both sides collapse; and every block includes
$r=1$, where $e_{r-1}=e_0=1$. Rick's numerical check covered $r=2,\dots,5$ only, at a single
$m$ large enough to be generic, so these two edges are new.

The transposition formula of Remark~\ref{rem:transp} was checked separately, against the same
brute-force $\sm$, on all four (L)-inputs at $m=4$ and on $x_1^a$, $a\le3$, at $m=2$.

Scripts: \texttt{proofs/scripts/day2026-09-18/verify\_fast.py} (the 270 checks above) and
\texttt{proofs/scripts/day2026-09-18/sigma\_bruteforce.py} (the first instrument, which uses
rational-function arithmetic and carries the Remark~\ref{rem:transp} check). Log:
\texttt{proofs/scripts/day2026-09-18/verify\_fast\_output.txt}.

\section{What this says about the handoff}

Rick's handoff suggested Macdonald III.5 (Pieri for Hall--Littlewood polynomials) or Ram--Yip
alcove walks. Neither is needed, and in hindsight neither is natural: the right-hand sides look
like a Pieri expansion --- the four partitions of $r+2$ obtained from $(r)$ by adding two boxes
in at most three rows --- but that shape is an artifact of the \emph{left}-hand sides, which are
$x_1^{a}$ times a tail-symmetric function of degree $r+2-a$ for $a=1,2,3$. The genuine content
is a single closed form, Theorem~\ref{thm:master}, in which the Hall--Littlewood input enters
only as the one-row family $\qq_n=P_{(n)}$ --- and only through $n\le a-1\le2$. The whole of
(L1)--(L4) uses just $\qq_1=e_1$ and $\qq_2=e_1^2-\qt2 e_2$.

The operative fact is Proposition~\ref{prop:sigmapower}: $\sm$ sends $x_1^{a}\mapsto P_{(a)}$.
Everything else is $\Lm$-linearity. Since $\{x_1^{a}\}_{a\ge0}$ generates the ring of
$x_2,\dots,x_m$-symmetric polynomials as a $\Lm$-module, Proposition~\ref{prop:sigmapower}
determines $\sm$ completely on the domain where it is being used; any further identity of this
family is now a finite computation with no new input.

\end{document}
